\title[What Does a Multimodal Engagement Decoder Learn? An Audit on NeuMa]{What Does a Multimodal EEG--Eye-Tracking Engagement Decoder Learn? A Leakage-Aware Audit on the NeuMa Dataset}

\author[1]{\fnm{Priyadharshini} \sur{D}\textsuperscript{[0009-0004-2291-1814]}}\email{priyadharshini.2024b@vitstudent.ac.in}
\author*[1]{\fnm{Shridevi} \sur{S}\textsuperscript{[0000-0002-0038-7212]}}\email{shridevi.s@vit.ac.in}

\affil*[1]{\orgdiv{Centre for Neuroinformatics}, \orgname{Vellore Institute of Technology}, \orgaddress{\city{Chennai}, \country{India}}}

\abstract{
Multimodal decoders combining electroencephalography (EEG) and eye tracking increasingly predict consumer engagement and are judged by accuracy against a rule-defined label. Accuracy alone does not say which information a decoder uses. We audit this question on the public NeuMa dataset (42 participants) with a rule-defined engagement index that we construct from the recordings, a full multimodal architecture as the object of study, and EEGNet, ShallowConvNet and an eye-tracking long short-term memory network (LSTM) as references, under leave-one-subject-out cross-validation. Four findings result. First, the gaze inputs account for the accuracy: a probe on the index's five gaze terms reaches an area under the receiver operating characteristic curve (ROC-AUC) of 0.92, the gaze-only LSTM 0.89, and no multimodal model of the principal comparison exceeded its ROC-AUC (audited decoder 0.88; EEGNet or ShallowConvNet fused with the LSTM 0.81--0.83). Second, no tested EEG encoder outperformed a linear probe on the index's five EEG terms (0.49--0.64 against 0.67), although the same recordings separate resting state from browsing at 0.85--0.94, and fold-wise relabelling left this unchanged. Third, performance alone did not substantiate the decoder's connectivity, efficiency or interpretability mechanisms: density-matched topology surrogates produced no detectable loss, the spiking front-end was slower than a dense twin, and the rule module's decisions were bypass-dominated. Fourth, on a product-selection label, session-level dwell probes reached 0.85--0.91, no 5-s epoch model exceeded 0.61, and none added detectably to dwell. The contribution is these findings on the interpretation of multimodal decoding performance, together with the eight-step audit, which is released as a protocol with the code.
}

\keywords{Consumer engagement, Neuromarketing, Electroencephalography, Eye tracking, Multimodal learning, Subject-independent evaluation, Label construction, Audit protocol}

\maketitle

\section{Introduction}\label{sec:intro}

Consumer engagement, the degree to which a consumer attends to, processes and is moved by a marketing stimulus, is linked to advertising effectiveness, brand recall and purchase intention \cite{ref1,ref2}, and self-report measures of it are retrospective, introspective and biased \cite{ref2,ref3}. Consumer neuroscience has therefore turned to physiological signals, with electroencephalography (EEG) dominant for its temporal resolution, portability and cost \cite{ref1,ref6,ref7}, and, increasingly, to multimodal decoders that pair EEG with eye tracking (ET), on the rationale that EEG captures processing without gaze and ET captures gaze without processing \cite{ref8,ref11}. Such decoders are now built from the components that dominate EEG deep learning, graph neural networks over functional connectivity \cite{ref19,ref20,ref21,ref22}, attention-based cross-modal fusion \cite{ref12,ref24,ref25,ref26,ref27} and, on the NeuMa data specifically, EEG graph signal processing combined with gaze dynamics \cite{ref18}. The practice that these studies share is to report the accuracy of the decoder against an engagement label and to read a high value as evidence that the neural signal carries the construct. A high accuracy does not, however, say which source of information the decoder picked up \cite{ref64}: the neural signal, the concurrent gaze, or the construction of the label, which in consumer neuroscience is itself typically a rule over physiological and behavioural statistics. The question this paper asks is therefore not whether a multimodal engagement decoder can be accurate, but what it learns when it is.

The question is acute on NeuMa \cite{ref35}, to our knowledge the only public dataset with synchronised EEG, eye tracking and consumer annotations acquired during a realistic brochure-browsing and purchase task. Its authors describe the resource as exploratory, ask that strict statistical significance tests accompany any claim drawn from it, validate it by signal quality rather than by decoding, and define the product-level segments of the recording by the time the gaze spent inside each product's bounding box \cite{ref35}; the release annotates product selection but no engagement, so the engagement index used here is a fixed rule over EEG band-power and gaze statistics that we constructed from the recordings (Section~\ref{sec:label}). Audits in neighbouring fields show what can go wrong when this is not taken into account: a celebrated EEG decoder was found to classify the temporal structure of its block design rather than the stimuli \cite{ref61}; subject-mixed splits inflate EEG accuracies from near-perfect to near-chance when corrected \cite{ref62,ref28,ref29,ref30,ref31,ref32}; cross-validation on cohorts of this size carries error bars of about ten accuracy points \cite{ref63}; and when a label is confounded, the source of a decoder's information is ambiguous unless the confound is controlled fold by fold \cite{ref64}. Consumer neuroscience has its own precedent: physiological measures were shown to add to traditional advertising measures only when their incremental validity was tested directly, modality by modality, on one footing \cite{ref65}. None of the multimodal engagement decoders cited above has been evaluated in this way.

We built one and audited it. The system under audit, INS-HDGS-CMT, combines the components the field currently favours: dynamic functional-connectivity graphs encoded by graph attention, a spiking front-end, a temporal transformer, cross-modal fusion with the gaze stream, subject-invariance regularisers and a differentiable soft-rule module (Section~\ref{sec:system}). It was trained under leave-one-subject-out cross-validation (LOSOCV) with every baseline tuned per fold under the same early-stopping rule, and no tested baseline exceeded it after correction (Section~\ref{res:rq1}); on a contrast the recordings are expected to carry, resting state versus browsing, it decodes at ROC-AUC 0.99 (Section~\ref{res:rq2}), and it runs in real time on a CPU without per-person calibration (Sections~\ref{sec:system} and~\ref{res:robust}). It is retained here as the object of study rather than as the proposal, for two reasons. An audit of a deliberately weak model would show nothing, whereas an audit of a model that performs comparably to the best tested baselines shows where that performance comes from. And its components are the ones whose contribution the literature assumes, so the controls that attribute its accuracy speak to the field's designs, not to one architecture. As reference encoders we use EEGNet \cite{ref50}, ShallowConvNet \cite{ref51} and an eye-tracking LSTM, each alone and the two convolutional encoders in late fusion with the LSTM, so that every finding is anchored on standard parts.

The study has four objectives, each stated as a question that a specific control answers:
\begin{enumerate}
\item[RQ1] \textbf{Where does the accuracy come from?} Whether a multimodal decoder of the NeuMa index draws on the neural signal, on gaze, or on the label's own construction, tested by quantifying the linear recoverability of the label from each modality's defining terms, by removing each gaze input from the decoder under paired folds, and by comparing the decoder with gaze-only and fusion references on one footing.
\item[RQ2] \textbf{Does any tested EEG encoder outperform a probe on the index's own EEG terms?} Tested by a gaze-free variant of the decoder and eight tuned EEG encoders against a linear probe on the five EEG terms, repeated under a fold-wise relabelling that removes the held-out participant from the construction of its own labels, and accompanied by a positive control, resting state versus browsing, on the same recordings.
\item[RQ3] \textbf{Do the components contribute what they were designed to?} Tested by ablation of each component under paired folds, by density-matched static and random graph nulls that change only the connectivity topology, by direct measurement of the spiking front-end's cost and signal, and by a fidelity test of the rule module's gate.
\item[RQ4] \textbf{Does anything inside a 5-s epoch predict product selection?} Tested by training the decoder and the reference encoders on a behavioural label, whether the viewed product was selected for purchase, and comparing them with within-epoch linear probes and with session-level dwell-time probes.
\end{enumerate}

Correspondingly, the paper makes four findings and one methodological contribution.
\begin{itemize}
\item \textbf{Finding 1, the gaze inputs account for the accuracy.} A probe on the index's five gaze terms reaches ROC-AUC 0.92 and one on its five EEG terms 0.67. The gaze-only LSTM reaches 0.89 (balanced accuracy 0.76), and no model of the principal comparison that sees gaze exceeded its ROC-AUC: the decoder reaches 0.88 (0.75) with its whole architecture, and EEGNet or ShallowConvNet fused with that LSTM 0.81--0.83 (0.72--0.73); no eye-tracking or fusion model differed from the decoder after correction, and removing the gaze inputs from the decoder lowers its ROC-AUC by 0.29 (Section~\ref{res:rq1}).
\item \textbf{Finding 2, no tested EEG encoder outperformed a probe on the index's EEG terms, although the recordings carry browsing state.} Nine EEG encoders, EEGNet, ShallowConvNet and the gaze-free decoder among them, decode the index at ROC-AUC 0.49--0.64, with no significant difference among them and none above the 0.67 of the linear probe; fold-wise relabelling of the retested models and of both label-term probes changes none of these results. On the same recordings the convolutional, linear and graph encoders separate resting state from browsing at 0.85--0.94, so the recordings and pipelines carry condition-discriminating information (Section~\ref{res:rq2}).
\item \textbf{Finding 3, performance alone did not substantiate the decoder's mechanisms.} Replacing the measured connectivity by density-matched static or random graphs produced no detectable reduction, so the ablation cost of the graph pathway (0.09 in balanced accuracy) is not attributable to the measured topology. The spiking front-end reached 0.55 alone on the positive control and ran slower than a dense twin on the hardware used. Replacing the cross-modal transformer by cross-attention produced no detectable difference, and the rule module's decisions were bypass-dominated; a variant trained with the gate closed decides through its rules at a cost of two to three points of balanced accuracy (Section~\ref{res:rq3}).
\item \textbf{Finding 4, session-level dwell predicts product selection, and nothing tested inside a 5-s epoch adds to it.} Dwell probes reach ROC-AUC 0.85--0.91 on whether a viewed product was selected for purchase. Models of the 5-s epoch reach 0.48--0.61, and adding the best of them to a dwell probe changes its ROC-AUC by 0.005 (95\% CI $-0.001$ to $+0.012$) (Section~\ref{res:rq4}).
\item \textbf{Protocol.} The eight controls above, label recoverability, gaze-free variant, reference encoders on one footing, fold-wise relabelling, topology nulls, rule fidelity, positive control and product-selection label, form an audit that depends on no property of this architecture and is released with the code (Section~\ref{sec:audit}).
\end{itemize}

The architecture is reported with the same completeness as its audit, because a multimodal decoder evaluated without such controls would be credited with mechanisms that the controls do not substantiate. Section~\ref{sec:methods} describes the dataset and the label, the system under audit, the reference encoders and the audit protocol. Section~\ref{sec:results} reports the results in the order of the four questions, followed by the attribution summaries of the rule-only variant and by the robustness checks. Section~\ref{sec:discussion} interprets the findings for consumer neuroscience and for the design of multimodal decoders, states the limitations that bound them, and Section~\ref{sec:conclusion} concludes.


